\section{Tail-Guided Auditing and Calibration}
\label{sec:methods}

\subsection{Spatial marks and the CT benchmark}

Each patient has a fixed candidate roster. Proposal $i$ has image features $X_i$ and weak marks from roles A and B, each recorded as marked, assessed and unmarked, or unavailable. The original reference pairs spatially concordant nodule contours from roles C and D. A proposal label $Y_i=1$ means it matches a target under the original geometric rule. An \emph{omission} is an original target with no link to A's marks. This definition precedes scoring and includes targets without an eligible proposal.

We use official LUNA16 V2 coordinates and discard the provided candidate class \cite{setio2017,luna16v2}. The roster merges CAD outputs, including components developed on the LIDC source data. We assign roles by a seeded per-scan permutation of the four final LIDC reading sessions. They are unblinded records with shared initial marks \cite{armato2011}. The original target geometry is the mean center and diameter of paired C/D contours, linked one-to-one within the greater of 3\,mm and the mean contour radius. The supplement specifies the matcher and the complete cohort. The final fit, calibration and test sets contain 445, 175 and 265 distinct patients, with 92 observed training omissions and 42 test omissions. We use official subsets 0, 1 and 3--5 for fitting, 2 and 6 for calibration, and 7--9 for testing.

\subsection{A probability audit for the score tail}

We freeze a development-only first-read image score $s_i$ and its low-alert cutoff $t_0$. For within-scan rank $r_i$ in descending score order, let
\begin{equation}
 h_i=\max\!\left\{\exp[-(r_i-1)/6],\
                    \sigma((s_i-t_0+1)/0.7)\right\}.
 \label{eq:leverage}
\end{equation}
The first term covers proposals that dominate their own scan. The second favors proposals near or above the common cutoff. Image-tail sampling uses $h_i$ alone. Paired tail also estimates the union of the weak marks out of fold using an image model. For its prediction $p_i$ and zero, one or two observed marks, the working probability $u_i$ is $0.35p_i$, $0.55+0.35p_i$, or $0.92+0.06p_i$, respectively. The resulting variance proxy $v_i=\max\{u_i(1-u_i),0.006\}$ gives extra weight to informative disagreements.

Given budget $B$, the design minimizes the surrogate $\sum_i \bar h_i^2v_i(1/\pi_i-1)$ with $\bar h_i=\max(h_i,0.002)$ under a sum constraint and box bounds. The solution is
\begin{equation}
 \pi_i=\operatorname{clip}(c\bar h_i\sqrt{v_i},\,0.005b,\,0.95),
 \qquad \sum_i\pi_i=B.
 \label{eq:allocation}
\end{equation}
Here $b$ is the budget factor and $c$ sets the expected number of candidate decisions. Image tail uses constant $v_i$. Mark stratification uses $(0.02,0.35,0.20)b$ for zero, one and two marks and sets the expected budget for every policy. Uniform selection uses $B/N$ for $N$ candidates. Inclusion is Bernoulli with known positive probabilities. Audit policies share random uniforms, and every training objective within a policy shares the same revealed labels.

For a target $j$ with eligible proposal indices $\mathcal I_j$, the exact chance of exposure is
\begin{equation}
 a_j=1-\prod_{i\in\mathcal I_j}(1-\pi_i).
 \label{eq:reach}
\end{equation}
The empty product is one. Summing over omissions counts distinct reached targets in expectation and separates acquisition from duplicate candidate decisions.

\subsection{Image-only training}

The fixed-roster factorial uses either 32 handcrafted intensity and Hessian descriptors \cite{li2003} or a 32-dimensional embedding from a 3D CNN on 24\,mm candidate patches. The encoder is trained on first-role marks from 267 development patients and frozen. All learned heads in the observed arm share the same embedding. A controlled bookmark arm retains A's largest contour and uses its own first-role encoder. Every deployed head takes image features alone.

Let $R_i$ indicate an audit and let $q_i$ be a reference prediction trained outside the patient's fold. The discriminative auxiliary uses the image embedding and weak-mark indicators, including agreement and availability. The final ridge logit is fit with label
\begin{equation}
 z_i=q_i+\frac{R_i}{\pi_i}(Y_i-q_i).
 \label{eq:augmentation}
\end{equation}
Conditional on the pre-audit inputs and the out-of-fold fit, $\mathbb E_R[z_i]=Y_i$. Logistic loss is affine in its target, so the fixed-head risk has the reference expectation. The signed residual variance is $(1-\pi_i)(Y_i-q_i)^2/\pi_i$. The supplement gives image-only and EM auxiliaries, power tuning, and the weighted-loss ablation. The factorial also includes first- and second-read fitting, nonnegative PU with an audited prior, direct inverse-probability-weighted (IPW) fitting, and a full-reference head.

\subsection{Patient alerts and reference membership}

A patient score is the maximum candidate logit $S_j$. For $n$ reference-negative calibration patients, use the ascending order statistic
\begin{equation}
 T_\alpha=S_{(k)},\qquad k=\lceil(n+1)(1-\alpha)\rceil ,
 \label{eq:cutoff}
\end{equation}
with an infinite cutoff if $k>n$. A proposal alerts on strict exceedance. With distinct exchangeable maxima from the stated negative population, exceedance probability is $(n+1-k)/(n+1)$ \cite{angelopoulos2023conformal}. A target is recovered if any of its proposals exceeds the cutoff. Duplicate alerts count toward FP/scan.

We report two analyses. The locked retrospective comparison uses the original C/D reference and nominal $\alpha=0.10$. Development selected half-budget paired-tail discriminative augmentation against full-budget mark IPW. The original primary endpoint is recall on observed omissions, with a five-point noninferiority margin and prespecified alert and FP/scan limits for the efficiency comparison. Intervals resample test patients jointly across six matched audit draws. The supplement separates audit variation from calibration variation and reports the calibration-only matched-alert check. Earlier work inspected regions of the public evaluation scans. This study froze the weights, endpoints and contrasts before revision scoring on the designated calibration and test subsets.

The score-fixed follow-up changes only which patients qualify as negative. Let $P$ be the C/D negatives and let $U$ be the patients negative under both C/D and the official LUNA16 reference, which requires agreement from at least three of four readers \cite{setio2017}. The $P$ score distribution mixes maxima from $U$ and from $P\setminus U$. Large maxima in the latter group can raise the upper quantile. We preserve all candidate scores and all original test targets, and evaluate both cutoffs on the same $U$-negative test patients. Thus the threshold effect on recovery is isolated from both score changes and changes to the target denominator.

We use 8,192 shared random deletions within calibration subsets, matching the number of exclusions in each subset. With $m$ removals the $k'$th remaining score originally lay between ranks $k'$ and $k'+m$. The random-deletion distribution measures the rank shift expected from the smaller pool. We also rebuild calibration pools for every pair of reader roles, preserving the test target list and using patients negative by every pair and the official reference for the common alert denominator. These follow-ups extend an initial anchor check on the same public cohort.
